% Generated draft: requires mathematical and source review.
\documentclass[11pt]{article}
\usepackage[margin=25mm]{geometry}
\usepackage{fontspec}
\setmainfont{Latin Modern Roman}
\usepackage{amsmath,amssymb,mathtools,unicode-math}
\setmathfont{Latin Modern Math}
\usepackage{enumitem,needspace,xurl,hyperref}
\setlength{\parindent}{0pt}
\setlength{\parskip}{5pt}
\setlength{\emergencystretch}{3em}
\newcommand{\N}{\mathbb{N}}
\newcommand{\Z}{\mathbb{Z}}
\newcommand{\Q}{\mathbb{Q}}
\newcommand{\R}{\mathbb{R}}
\newcommand{\C}{\mathbb{C}}
\newcommand{\F}{\mathbb{F}}
\newcommand{\A}{\mathbb{A}}
\newcommand{\PP}{\mathbb{P}}
\newcommand{\E}{\mathbb{E}}
\newcommand{\eps}{\varepsilon}
\newcommand{\dd}{\,\mathrm{d}}
\begin{document}
\section*{UnsolvedMath ID: 30005300}
\textbf{Source problem code:} catalogue entry
\subsection{Definitions and mathematical statement}
The supplied question asks for a characterization of matroids admitting complete-intersection realizations by points. It does not specify the ground field, its characteristic, the ambient projective space, or whether complete intersection is meant scheme-theoretically. Consequently, the following gives a precise parameterized interpretation, whose conventions require confirmation against the report's surrounding discussion.

Fix a field \(k\). Let \(M=(E,\rho)\) be a finite simple matroid of rank \(r\geq 1\), where \(\rho:2^E\to\Z_{\geq 0}\) is its rank function. Here simplicity means that \(M\) has no loops or parallel elements. A realization by distinct \(k\)-rational projective points is a family
\[
Z=\{p_e=[v_e]:e\in E\}\subseteq\PP^{r-1}_k,
\qquad v_e\in k^r\setminus\{0\},
\]
with distinct \(p_e\), satisfying
\[
\rho(A)=\dim_k\operatorname{span}_k\{v_e:e\in A\}
\quad\text{for every }A\subseteq E.
\]
In particular, these points span \(\PP^{r-1}_k\). Give \(Z\) its reduced closed subscheme structure and let \(I(Z)\subseteq S=k[x_0,\ldots,x_{r-1}]\) be its homogeneous vanishing ideal. In the scheme-theoretic interpretation, \(Z\) is a complete intersection if
\[
I(Z)=(f_1,\ldots,f_{r-1})
\]
for a homogeneous regular sequence of positive-degree forms. Regularity means that multiplication by \(f_i\) is injective on \(S/(f_1,\ldots,f_{i-1})\) for each \(i\), and the generated ideal is proper. For \(r=1\), this is the empty regular sequence, with ideal \((0)\), giving the single reduced point of \(\PP^0_k\).

Under these conventions, the mathematical task is: for a specified field \(k\), characterize the finite simple matroids \(M\) for which such a realization \(Z\) exists and is a complete intersection. The source additionally asks for necessary combinatorial conditions beyond unspecified ``obvious'' conditions; that phrase cannot be made more precise from the supplied excerpt.

The restrictions to distinct rational points and to a reduced scheme-theoretic complete intersection are explicit interpretive choices, not hypotheses established by the supplied record. The intended field quantifier, treatment of nonsimple matroids, empty configurations, and rank-zero cases remain unresolved. The source mentions Terao's question about whether arrangement freeness is combinatorial as motivation; it is not identified with the characterization problem above.
\subsection{Short English statement}
Characterize the matroids that can be represented by projective points which themselves form a complete intersection. A precise version asks, over a specified field, for distinct spanning points whose linear-dependence matroid is the given matroid and whose reduced homogeneous ideal is generated by a regular sequence. The supplied excerpt leaves the field and complete-intersection conventions unspecified.
\subsection{Sources}
\begin{itemize}
\item Supplied local UnsolvedMath registry record, permanent ID \(30005300\), title \emph{Complete-Intersection Realizations of Matroids}, exact source code \texttt{catalogue entry}; registry locations: \texttt{data/unsolvedmath/problems.json} and the supplementary \texttt{catalog/problems.json}. These are provenance references for the supplied record, not a claim of independent file inspection.
\item Supplied source citation: \emph{Mini-Workshop: Subvarieties in Projective Spaces and Their Projections} (2023), Oberwolfach Reports, pp.\ 3093--3119, DOI \texttt{10.4171/owr/2022/54}. The extracted passage carries page number 3107. \url{https://doi.org/10.4171/owr/2022/54}. The record gives proposed year 2022 and names no proposer.
\item Supplied literature reference: Abe, \emph{Combinatorial data of a free arrangement and Terao's conjecture}. \url{https://arxiv.org/abs/1307.3750}.
\item Supplied literature reference: EMS Press, OWR 2022/54. \url{https://doi.org/10.4171/owr/2022/54}.
\item Literature assessment dated exactly \texttt{2026-08-22}: the supplied curation labels the record partially solved and states that Terao's freeness conjecture remains open in general, with results for important classes and finite ranges and substantial computational verification; it also states that the extracted motivating question is unsettled. This is supplied dated curation, not an independently verified literature assessment, and it does not identify partial results for the complete-intersection characterization itself.
\item The supplied registry labels the statement \texttt{corrected\_verified} and reports comparison with the official report. That verification claim is preserved as source metadata; no source visit or independent mathematical verification is asserted here.
\end{itemize}
\end{document}
